\subsection{周期函数}

定义 1. 函数 \(f\) （定义在 \(\mathbf{R}^n\) 上、取值于任意集合 \(\mathbf{E}\) 中的）称为周期的，如果存在一点 \(a \neq 0\) （ \(\mathbf{R}^n\) 中的），使得

\[
f (\boldsymbol {x} + \boldsymbol {a}) = f (\boldsymbol {x})\tag{1}
\]

对一切 \(x \in R^{n}\) 成立。若 f 是周期的，则使关系 (1) 成为关于 x 的恒等式的每个点 \(a \in R^{n}\) 都称为 f 的一个周期。

周期函数 f 的所有周期所成的集合 G 显然是加法群 \(R^{n}\) 的一个子群（由假设它不只由 o 组成）。若 f 是 \(R^{n}\) 到豪斯多夫拓扑空间 E 中的连续周期映射，则它的周期群 G 是闭的。因为若用 \(G_{x}\) 表示由一切 \(a \in R^{n}\) （使得 \(f(x + a) = f(x)\) ，对给定的点 \(x \in R^{n}\) ）组成的集合，则 G 是诸 \(G_{x}\) （当 x 取遍 \(R^{n}\) 时）的交，而每个 \(G_{x}\) 都是闭的（第 I 章 § 8 第 1 小节命题 2）。设 V 是含于 G 中的最大向量子空间（第 2 小节定理 2）；函数 f 在模 V 的每个陪集上都取常值；若 W 是与 V 互补的向量子空间，则 f 由它在 W 上的限制确定。换言之（W 是同构于某个 \(R^{p}\) 的拓扑群，p 是某个整数）， \(R^{n}\) 上连续周期函数的研究可归结为这类函数中周期群 G 为离散群者的研究；若这个群的秩为 q，则称函数 f 是 q 重周期的，而生成 G 的 q 个点的每个自由系都称为 f 的一个基本周期系。

若 \((\boldsymbol{a}_{i})\) 与 \((\boldsymbol{b}_{i})\) 是 f 的两个基本周期系，则我们已看到（第 1 小节），每个都可以由另一个通过一个以整数为系数、行列式为 \(\pm1\) 的线性变换得到。

设 \(\varphi\) 是 \(\mathbf{R}^n\) 到 \(\mathbf{R}^n/\mathbf{G}\) 上的典范映射；对每个映射 \(g\) （ \(\mathbf{R}^n/\mathbf{G}\) 到集合 \(\mathbf{E}\) 中的），对应着函数 \(\dot{g}=g\circ\varphi\) ，它是 \(\mathbf{R}^n\) 到 \(\mathbf{E}\) 中的周期映射，其周期群包含 G；反过来，每个周期群包含 G 的 \(R^{n}\) 到 E 中的映射都具有这种形式，因为它与关系 \(x \equiv y \pmod{G}\) 相容（《集合论》R， § 5 第 7 小节）。这样我们就定义了一个双射 \(g \to \dot{g}\) ，它把 \(R^{n}/G\) 到 E 中的一切映射所成的集合映到 \(R^{n}\) 到 E 中且周期群包含 G 的一切映射所成的集合上。为使 \(\dot{g}\) 连续（E 是拓扑空间），必要充分条件是 g 连续（第 I 章 § 3 第 4 小节命题 6）。

